Практическое применение большой языковой модели связано с этапом вывода: модель получает пользовательский запрос, выполняет прямые проходы с уже зафиксированными параметрами и возвращает результат. Такой этап работы модели принято называть инференсом. Зачастую повышение качества больших языковых моделей на широком наборе задач достигается за счёт увеличения количества их параметров и вычислительного бюджета~\cite{kaplan2020scaling,hoffmann2022training}. Однако в традиционных плотных архитектурах трансформеров все параметры соответствующего слоя участвуют в вычислениях на каждом проходе модели, поэтому масштабирование таких моделей приводит к быстрому росту стоимости инференса~\cite{vaswani2017attention,brown2020language}.

Эта проблема важна не только для отдельных моделей. Современные SOTA модели уже достигают масштаба в сотни миллиардов параметров: DeepSeek-V3.2 содержит 671 миллиардов параметров~\cite{deepseek2025v32}, а Switch Transformer и GLaM демонстрируют масштабирование MoE архитектур до триллионнов параметров~\cite{fedus2021switch,du2021glam}. Для таких моделей веса, KV кэш и промежуточные активации становятся слишком большими для одного GPU. Поэтому распределенный инференс является необходимым условием практического применения крупных моделей~\cite{aminabadi2022deepspeed}.

Одним из методов распределенного инференса является конвейерный параллелизм (pipeline parallelism, PP), при котором модель разбивается на последовательные части, размещаемые на разных GPU. Однако данный подход имеет ограничения для инференса больших языковых моделей из-за последовательного характера выполнения вычислений между стадиями конвейера, что приводит к увеличению времени выполнения отдельных запросов и неполной утилизации вычислительных ресурсов. Другим широко используемым подходом является тензорный параллелизм (tensor parallelism, TP), при котором вычисления внутри отдельных операций, например матричных умножений, распределяются между несколькими устройствами~\cite{shoeybi2019megatron,narayanan2021megatron}. В отличие от PP, TP позволяет сохранять низкое время выполнения для отдельных запросов и более эффективно загружать GPU, однако современные работы по тензорному параллелизму при инференсе отмечают, что по мере роста больших языковых моделей и распределения модели на большее число устройств коммуникационная стоимость увеличивается и может становиться узким местом производительности~\cite{dong2024lowbit}. В результате даже сочетание различных методов параллелизма не позволяет масштабировать плотные архитектуры без существенного роста стоимости инференса.

Поэтому была предложена альтернативная архитектура: Mixture-of-Experts (MoE)~\cite{shazeer2017moe}. Суть данного подхода заключается в том, что в блоке трансформера полносвязный линейный слой (feed-forward network) заменяется на множество меньших полносвязных блоков (экспертов), а также добавляется роутер, который определяет, какие из экспертов будут активны на текущем прямом проходе.

Ключевое преимущество MoE заключается не в том, что на каждом проходе используется фиксированная доля всех весов модели, а в том, что активируется только часть экспертных параметров. Плотные части модели продолжают выполняться полностью, однако вычисления в экспертных слоях становятся разреженными. За счет этого MoE позволяет увеличивать общую емкость модели без пропорционального роста стоимости одного прохода~\cite{fedus2021switch,du2021glam}.

Механизм масштабирования MoE связан с экспертным параллелизмом (expert parallelism, EP): разные эксперты размещаются на разных видеокартах, а токены после роутинга отправляются к выбранным экспертам. В классической схеме для этого используется примитив коллективных коммуникаций All-to-All: сначала активации передаются от одного ранга ко всем остальным, и так для каждого, затем результаты экспертных вычислений собираются обратно~\cite{rajbhandari2022deepspeedmoe,deepseekDeepEP}. Такие коллективные операции на NVIDIA GPU обычно реализуются поверх специализированных коммуникационных библиотек, в частности NCCL и NVSHMEM~\cite{nvidiaNCCLDocs,nvidiaNvshmem}.

Несмотря на высокую производительность, коллективная схема плохо подходит для динамической балансировки нагрузки. Роутер MoE выбирает экспертов в зависимости от потока входных данных, поэтому часть экспертов может регулярно активироваться на гораздо большем количестве токенов, чем остальные. На коллективах раскладка экспертов по рангам фиксируются заранее, а каждый шаг All-to-All связывает всех участников. Поэтому система не может во время работы свободно перенаправлять запросы к менее загруженной реплике эксперта, исключать перегруженный экспертный блок или добавлять новый без полного перезапуска всей системы.

Таким образом, возникает вопрос о том, можно ли заменить статичную коллективную схему взаимодействия в экспертном слое на архитектуру, позволяющую принимать решения о выборе экспертного блока во время работы системы.

\subsection*{Цели и задачи}
\addcontentsline{toc}{subsection}{Цели и задачи}

Основной целью данной работы является предложить и разработать рабочий прототип альтернативного подхода к инференсу больших языковых моделей, основанных на архитектуре MoE, удовлетворяющий требованиям динамической балансировки нагрузки и обеспечивающий более высокую отказоустойчивость по сравнению с оригинальным подходом, основанным на коллективах.

Для достижения поставленной цели была выдвинута следующая гипотеза: замена коллективных All-to-All коммуникаций на уровне экспертного слоя в MoE моделях на клиент-серверное взаимодействие между основной частью модели и экспертными блоками позволяет выбирать конкретного исполнителя во время работы системы, перераспределять нагрузку между репликами экспертов и сохранять доступность при отказе части экспертных блоков.

Для проверки данной гипотезы ставятся следующие задачи:

\begin{itemize}
    \item Разработать интерфейс для высокопроизводительной передачи данных между GPU, а также две реализации этого интерфейса: для NIXL (NVIDIA Inference Xfer Library) и MoonCake Transfer Engine
    \item Разработать клиент-серверную архитектуру для распределенного инференса MoE моделей с поддержкой асинхронного взаимодействия между основной частью модели и экспертными блоками
    \item Разработать систему асинхронной диспетчеризации запросов с поддержкой очередей
    \item Разработать механизмы динамической балансировки нагрузки между экспертными блоками, включая выбор менее загруженной реплики эксперта
    \item Реализовать систему мониторинга состояния экспертных блоков и механизмы обнаружения отказов в режиме реального времени
    \item Провести оптимизацию сетевого взаимодействия для минимизации накладных расходов при передаче активаций и результатов вычислений
    \item Провести оптимизацию GPU-вычислений для экспертных блоков
    \item Провести экспериментальное сравнение производительности предложенного подхода с существующими коллективными решениями
    \item Проверить возможность динамической балансировки нагрузки и сохранения доступности при частичных отказах
\end{itemize}

\subsection*{Достигнутые результаты}
\addcontentsline{toc}{subsection}{Достигнутые результаты}

В рамках данной работы был предложен и реализован прототип отказоустойчивого инференса MoE моделей, основанный на клиент-серверной архитектуре, подтверждающий выдвинутую гипотезу.

Для этого была разработан программный модуль, включающий следующие компоненты:

\begin{itemize}
    \item Интерфейс для высокопроизводительных коммуникаций
    \item Асинхронная клиент-серверная архитектура, предоставляющая возможность динамического масштабирование экспертных блоков
    \item Система мониторинга и обнаружения отказов экспертных блоков
    \item Механизмы динамической балансировки нагрузки между репликами экспертов
    \item Оптимизации GPU-вычислений
\end{itemize}

Экспериментальные оценки показали, что предложенный подход обеспечивает возможность динамической балансировки нагрузки между экспертными блоками и сохраняет доступность при их частичных отказах. Полученные результаты демонстрируют практическую применимость клиент-серверного подхода для построения гибких систем инференса MoE моделей.

\subsection*{Структура работы}
\addcontentsline{toc}{subsection}{Структура работы}

Работа состоит из введения, трех глав и заключения.

В первой главе приводится обзор литературы и существующих подходов к инференсу больших языковых моделей. В ней рассматриваются архитектура трансформера, особенности масштабирования плотных моделей, конвейерный, тензорный и экспертный параллелизм, архитектура Mixture-of-Experts, коллективная операция All-to-All, а также ограничения коллективных решений с точки зрения балансировки нагрузки и обработки частичных отказов.

Во второй главе описывается предложенный метод: клиент-серверная архитектура инференса MoE моделей, устройство программных компонентов MoE инференса и подсистемы передачи тензоров, схема взаимодействия основной части модели с экспертными блоками, механизм асинхронной диспетчеризации, мониторинг состояния экспертных блоков и динамическая балансировка нагрузки. Также в главе рассматриваются оптимизации, выполненные в ходе реализации: уменьшение объема передаваемых по сети данных, оптимизация передачи тензоров, выравнивание экспертных активаций и применение групповых матричных умножений. Архитектура проекта и основные компоненты метода иллюстрируются схемами.

В третьей главе приводятся эксперименты, направленные на оценку производительности предложенного подхода, стоимости передачи тензоров, эффективности оптимизаций, возможности динамической балансировки нагрузки и поведения системы при отказах экспертных блоков. В ней фиксируется компромисс между снижением абсолютной производительности и повышением гибкости управления экспертными блоками.
